\def\paperid{A01}
\documentclass[10pt,a4paper]{article}
\usepackage[margin=23mm,top=24mm,bottom=24mm]{geometry}
\usepackage{fontspec,polyglossia,amsmath,amssymb,booktabs,array,enumitem,xcolor,fancyhdr,xurl,hyperref,microtype}
\setdefaultlanguage{spanish}\setmainfont{Latin Modern Roman}\setsansfont{DejaVu Sans}
\definecolor{ink}{HTML}{17375C}\definecolor{mint}{HTML}{0F8F86}
\hypersetup{colorlinks=true,urlcolor=mint,linkcolor=ink,pdfauthor={Juan Manuel Goenaga Castro}}
\setlength{\parindent}{0pt}\setlength{\parskip}{5pt}\linespread{1.05}\emergencystretch=2em
\setlength{\headheight}{15pt}\pagestyle{fancy}\fancyhf{}
\fancyhead[L]{\sffamily\footnotesize GoToGym / Quantum Coach}\fancyhead[R]{\sffamily\footnotesize Robustecimiento QAF / \paperid}
\fancyfoot[L]{\scriptsize Prepublicación metodológica / 22 de octubre de 2026}\fancyfoot[R]{\small\thepage}
\setlist{nosep,leftmargin=1.5em}
\begin{document}
\begin{center}{\sffamily\small\textcolor{mint}{SERIE DE ROBUSTECIMIENTO / A01}}\par\vspace{10pt}{\sffamily\LARGE\bfseries\textcolor{ink}{Prosodia continua y ritmo narrativo}}\par\vspace{7pt}{\large Control temporal QAF sin silencios de relleno}\par\vspace{12pt}\textbf{Juan Manuel Goenaga Castro}\par GoToGym Quantum Research\par 22 de octubre de 2026 / versión 1.1\end{center}\colorbox{ink!6}{\parbox{.96\linewidth}{\small\textbf{Estado.} Prepublicación metodológica sin revisión por pares. Propuesta computacional implementada; evaluación perceptual y estudio con participantes pendientes.}}\begin{abstract}Se propone separar tiempo de película, tiempo de habla y tiempo de lectura visual. La contribución implementada es un selector restringido de velocidad, conservación de los descansos originales y medición energética explícita. El sistema conserva la fuente, evita cortar palabras y entrega una grabación continua más rápida que el montaje anterior. La identidad de voz, el acento y la naturalidad permanecen dimensiones distintas de la tasa léxica. Se presenta el problema observado, una formalización con unidades y restricciones, el código correspondiente y sus controles. Se distingue la verificación de ingeniería de calidad percibida, identidad regional y efecto comercial. El uso QAF corresponde a la región clásica identificable y conserva un comparador con los mismos recursos.\end{abstract}
\section{Diagnóstico y pregunta de investigación}
El episodio antecedente reducía el tempo al 82 por ciento y añadía once pausas para ocupar un minuto. El usuario identificó una voz lenta y descansos robóticos. Ese comentario es un hallazgo cualitativo individual, no una estimación de preferencia poblacional. El fallo de diseño consistía en imponer a la narración la duración de la película. Una cuenta de palabras puede orientar el guion, pero el presupuesto definitivo debe medirse sobre la señal real.

Se propone separar tiempo de película, tiempo de habla y tiempo de lectura visual. La contribución implementada es un selector restringido de velocidad, conservación de los descansos originales y medición energética explícita. El sistema conserva la fuente, evita cortar palabras y entrega una grabación continua más rápida que el montaje anterior. La identidad de voz, el acento y la naturalidad permanecen dimensiones distintas de la tasa léxica.
\section{Modelo matemático y restricciones}
N es la cantidad de palabras del guion; D la duración de la grabación continua, en segundos; r el factor de tempo; l la entrada de voz; T el presupuesto de película y w* la tasa objetivo, en palabras por minuto. El intervalo de tempo y la reserva final son parámetros editoriales versionados.
\begin{equation}w(r)=60Nr/D,\quad r\in[0.90,1.08],\quad D/r+l\leq T-0.4.\end{equation}

\begin{equation}L_r=\left(\frac{w(r)-w^*}{w^*}\right)^2+0.002\log^2r,\qquad \Delta p_j=0.\end{equation}

\begin{equation}e_k=\sqrt{\frac1M\sum_{n=kM}^{(k+1)M-1}y[n]^2},\quad q_k=\mathbf1[e_k<10^{-42/20}],\quad M=0.02f_s.\end{equation}
El término cuadrático compara tasas sin unidades y el logaritmo regulariza cambios de tempo. La condición de factibilidad protege la última palabra; la condición sobre las pausas prohíbe comprar duración con silencio insertado. La rejilla de candidatos incluye el valor que aproximaría la tasa objetivo, dentro del intervalo permitido. La detección de baja energía se realiza en ventanas de veinte milisegundos y excluye extremos. Respiraciones y consonantes sordas pueden quedar bajo el umbral: el resultado es un proxy acústico, no una segmentación lingüística.
\subsection{Región identificable QAF y comparador}
La integración usa el caso diagonal del formalismo, con base fija de candidatos y compuerta de factibilidad g:
\begin{equation}
p_j=\frac{g_j\exp[-(L_j-L_{\min})/\tau]}{\sum_k g_k\exp[-(L_k-L_{\min})/\tau]},\quad\rho=\operatorname{diag}(p),\quad\rho\succeq0,\quad\operatorname{Tr}\rho=1.
\end{equation}
Estos pesos describen decisión, no satisfacción del público. No se aprende ni asigna fase sin observaciones que la identifiquen. La enumeración nativa selecciona el mínimo factible; el comparador clásico recibe exactamente los mismos candidatos y coincide por construcción. No se reclama ventaja cuántica. MPS y TDVP siguen disponibles como referencias matemáticas anteriores, pero no intervienen en la síntesis de voz, el transporte de luz o una supuesta dinámica cerebral.

Con pérdidas desplazadas no negativas, se conserva la compilación triangular de selección única:
\begin{equation}
E(x)=\sum_j(\widetilde L_j-A)x_j+2A\sum_{j<k}x_jx_k+A,\qquad A>\max_j\widetilde L_j.
\end{equation}
La constante distingue correctamente el conjunto vacío; la interacción aparece una sola vez por pareja. La máscara de alternativas prohibidas es una restricción nativa, no una penalización blanda. Una prueba independiente enumera las 64 configuraciones de seis variables y compara con el óptimo restringido. La positividad y la traza se verifican por separado. Este uso explícito de la región clásica evita forzar complejidad tensorial sin evidencia incremental.
\section{Implementación y trazabilidad}
voice\_\allowbreak{}rhythm.py implementa natural\_\allowbreak{}voice y silence\_\allowbreak{}metrics; quality.py conserva alternativas, pérdidas y máscara factible. La transformación atempo evita cambio intencional de altura y se normaliza después el premáster. El worker vuelve a transcribir cuando se cambia el tempo. El servidor genera con presupuesto de episodio, sin exigir que toda voz dure sesenta segundos. La ruta antigua de ajuste fijo se mantiene por compatibilidad y no gobierna el nuevo episodio.

Módulos: \texttt{voice\_\allowbreak{}rhythm.py}, \texttt{quality.py}, \texttt{local\_\allowbreak{}voice\_\allowbreak{}worker.py}. Cada transformación conserva archivos de entrada, parámetros y hashes. No se sustituye una capacidad ausente por una declaración de disponibilidad.
\subsection{Controles que pueden refutar la implementación}
\begin{flushleft}\small test\_\allowbreak{}continuous\_\allowbreak{}timing\_\allowbreak{}never\_\allowbreak{}inserts\_\allowbreak{}sentence\_\allowbreak{}pauses\_\allowbreak{}or\_\allowbreak{}pads\_\allowbreak{}minute\\ test\_\allowbreak{}energy\_\allowbreak{}detector\_\allowbreak{}distinguishes\_\allowbreak{}edge\_\allowbreak{}silence\_\allowbreak{}and\_\allowbreak{}internal\_\allowbreak{}gap.\end{flushleft}

La batería combina invariantes algebraicos, casos adversos y render pequeño con decodificación. El conjunto de tests es un piso de reproducibilidad: aprobar una identidad de bytes no acredita veracidad de contenido, y una ecuación correcta no acredita preferencia artística.
\section{Resultados computacionales observados}
El factor pasó de 0,82 a 0.921951, un aumento relativo de velocidad de 12.43 por ciento. Las once pausas insertadas del relato pasan a cero. La narración completa mide 42.767 segundos; su mayor tramo interno de baja energía mide 0.54 segundos. El resto de la película se reserva a imagen y música.

La suite ejecutada contiene 72 pruebas aprobadas. Son contratos de software e integración, no 72 observaciones independientes de calidad artística. Los datos completos se registran en el informe de versión.

El archivo \texttt{docs/validation/robustness-v0.4.json} conserva medidas y hashes de entrega. El proyecto completo con voz se distribuye por separado del código público. No se publican notas originales, perfiles biométricos ni sus transcripciones privadas.
\section{Protocolo de evaluación y ablación}
Comparar tres versiones del mismo guion y referencia: montaje antecedente, voz continua sin cambio de tempo y voz continua con tasa seleccionada. Equilibrar nivel, orden y dispositivo; ocultar el nombre del método. Medir preferencia pareada, palabras omitidas y naturalidad ordinal con oyentes familiarizados con Barranquilla y otros oyentes. Rechazar el candidato si su mayor velocidad reduce comprensión o identidad. Estimar variabilidad con un piloto antes de calcular potencia confirmatoria.

La aceptación es jerárquica: primero se satisfacen las restricciones duras; después se optimiza calidad dentro del conjunto admisible. Se registrarán candidatos fallidos, tiempo y recursos, no solo la toma ganadora. El episodio, la persona y la familiaridad se controlan por separado. Repetir valoraciones de un único video no crea videos independientes. La muestra confirmatoria se calculará con variabilidad del piloto, sin inventar potencia o resultados.
\section{Discusión, límites y propiedad intelectual}
Una tasa mayor no garantiza naturalidad. atempo puede introducir artefactos, aun sin cambio intencional de tono. El detector no recorta respiraciones ni afirma detectar todas las pausas. El caso usa una única voz sintética y un episodio; no certifica acento costeño ni preferencia general.

Los aportes de integración se preparan a nombre de Juan Manuel Goenaga Castro. Se conserva atribución de Antoni Bancells al antecedente técnico QAF y de los modelos, métodos y fuentes de terceros. No se reivindica entrenamiento propio de Qwen o Whisper. Una ampliación relacional deberá superar un comparador clásico con las mismas entradas y presupuesto; si no aporta valor, la solución sencilla continúa siendo válida.
\section{Conclusión}
La propuesta computacional descrita se encuentra en el código entregado y permite repetir sus controles. La calidad percibida, la autenticidad vocal y la eficacia comercial requieren sus propios datos; no se declaran resueltas por completar una exportación.
\begin{thebibliography}{9}
\bibitem{qaf} Goenaga Castro, J. M. (2025--2026). Quantum Analysis Frameworks y serie QAF audiovisual I--V, documentos aportados. Bancells, A. (2026). Dinámica operacional, MPS/TDVP y compilación QUBO, antecedente técnico aportado. Atribuciones conservadas; sus bancos completos no se repitieron en esta ampliación.
\bibitem{ffmpeg} FFmpeg. Documentación primaria de filtros atempo, loudnorm, sidechaincompress y scale. Consulta 22/10/2026. \url{https://ffmpeg.org/ffmpeg-filters.html}.
\bibitem{qwen} Qwen Team (2026). Qwen3-TTS, tarjeta del modelo 12Hz-0.6B-Base. Modelo preentrenado de terceros. \url{https://huggingface.co/Qwen/Qwen3-TTS-12Hz-0.6B-Base}.
\end{thebibliography}
\end{document}
